\begin{proposition}[No pointwise physical-symmetry stabilizer for a regular grazing face]
\label{prop:no_pointwise_symmetry_stabilizer_grazing_face}
Let a finite group of isometries of the flat torus preserve a corner-free
dispersing table with pairwise disjoint strictly convex scatterers, and let it
act on the outgoing collision section by
\((q,v)\mapsto(Sq,DS\,v)\).  No nonidentity element fixes pointwise an open
subarc of a regular one-step physical grazing face.  Consequently, after
removing the finite union of closed fixed sets and shrinking inside a
nonempty complementary subarc, every finite geometric symmetry orbit of a
regular physical grazing face is free and has pairwise disjoint
representatives.  The preceding symmetry-restricted split
surjection therefore
applies.  In particular finite physical symmetry cannot supply the exceptional
pointwise stabilizer Ward identity left open by the abstract isotypic
calculation.
\end{proposition}

\begin{proof}
Work in a local Euclidean lift of the torus.  A nontrivial translation has no
fixed configuration point.  A nonreflection orthogonal isometry has only
isolated fixed configuration points and no open fixed set of phase states.
The only remaining local possibility is a reflection across a geodesic axis
\(A\).  A fixed phase state must have \(q\in A\) and \(v\) tangent to \(A\).

Suppose such states formed an open subarc of a one-step grazing face.  By
equivariance and uniqueness of the regular free flight, the corresponding
grazing impact point \(p\) would also be fixed by the reflection.  The target
scatterer containing \(p\) cannot be exchanged with a distinct scatterer:
then \(p\) would belong to two components, contrary to their positive
separation.  Hence that scatterer is reflection invariant.  A strictly convex
reflection-invariant boundary meets the reflection axis orthogonally; at each
intersection its normal is tangent to \(A\).  Since the flight velocity is
tangent to \(A\), the impact at \(p\) is normal, not grazing, a contradiction.

Thus every nonidentity group element has fixed set with empty interior in a
regular grazing face.  The group is finite, so the union of these fixed sets
still has empty interior.  Choose a smaller compact subarc away from it; its
stabilizer is trivial, and a further shrinking makes its finitely many images
pairwise disjoint.  The last assertions follow from the isotypic right-inverse
construction preceding this proposition.
\end{proof}

\begin{proposition}[Periodic-orbit rigidity of a local Ward primitive]
\label{prop:periodic_orbit_rigidity_local_ward_primitive}
Work at the base table, where \(\rho_0\equiv1\) relative to the invariant
collision flux and
\(\Ltransfer_0h=h\circ\Pcal_0^{-1}\) on the smooth invertible core.  Let
\(a\) be a centered branchwise \(C^r\) collision observable.  If its Poisson
equation has a single-valued continuous local primitive on the regular set,
\begin{equation}
\label{eq:local_ward_primitive_equation}
 a=(I-\Ltransfer_0)h,
\end{equation}
then every regular period-\(p\) point \(z\) satisfies
\begin{equation}
\label{eq:periodic_orbit_ward_condition}
 \sum_{k=0}^{p-1}a(\Pcal_0^kz)=0.
\end{equation}

For any \(N\) pairwise disjoint regular periodic orbits
\(\mathcal O_1,\ldots,\mathcal O_N\), the orbit-sum map on the centered smooth
observable core,
\begin{equation}
\label{eq:periodic_orbit_evaluation_map}
 \mathcal E_Na
 :=\left(\sum_{z\in\mathcal O_j}a(z)\right)_{j=1}^N,
\end{equation}
is onto \(\mathbb R^N\).  Hence observables admitting such a local primitive
lie in a closed subspace of codimension at least \(N\).  On a mixing
finite-horizon dispersing component, where arbitrarily large finite families
of disjoint regular periodic orbits are standard
\cite{Sinai1970,chernov2006chaotic}, this local-primitive class
has infinite codimension independently of the moving-face gluing obstruction.

Finally, a continuous local charge \(H\) satisfying
\(H\circ\Pcal_0=H\) is constant on the period-one ergodic component, so its
one-collision current is zero.  If instead one prescribes the increment
\(H-H\circ\Pcal_0^{-1}\), that source is exactly the observable-current class
of \cref{thm:natural_observable_current_recovery}.  Thus a local Noether charge
produces either the zero source or the already separated exact-current
benchmark; it does not supply the missing independently prescribed
nonconjugate correlation source.

There is also no missing finite-memory label rule.  Let \(\mathsf G\) be a
finite strongly connected directed higher-block graph coding a regular
collision subsystem, and let \(a\) depend only on one edge of \(\mathsf G\)
(equivalently, on a fixed finite collision word before passing to the
higher-block presentation).  If \eqref{eq:periodic_orbit_ward_condition}
holds for every periodic word in this subsystem, then there is a vertex
potential \(K\) such that, along the coded collision dynamics,
\begin{equation}
\label{eq:finite_memory_label_coboundary}
 a=K-K\circ\Pcal_0^{-1}.
\end{equation}
Thus every finite-label finite-memory local Ward rule that passes all coded
periodic tests is already an exact-current benchmark on that subsystem.
\end{proposition}

\begin{proof}
Evaluating \eqref{eq:local_ward_primitive_equation} along a regular periodic
orbit gives
\[
 a(\Pcal_0^kz)=h(\Pcal_0^kz)-h(\Pcal_0^{k-1}z).
\]
Summing over one period proves \eqref{eq:periodic_orbit_ward_condition}.

To prove surjectivity, choose one point \(z_j\in\mathcal O_j\) on each orbit
and pairwise disjoint small regular neighbourhoods meeting no other selected
orbit point.  Let \(b_j\) be a smooth bump supported in the \(j\)-th
neighbourhood with \(b_j(z_j)=1\).  Choose a smooth bump \(c\), supported away
from all selected orbits, with \(\int c\,d\nu_0=1\), and put
\[
 \widetilde b_j=b_j-\left(\int b_j\,d\nu_0\right)c.
\]
Then \(\widetilde b_j\) is centered and its orbit sums are the \(j\)-th
standard basis vector.  This proves that \(\mathcal E_N\) is onto.  Point
evaluation is continuous in the stated smooth observable topology, so the
common kernel is closed and has codimension \(N\).  The infinite-codimension
statement follows by taking arbitrarily large finite orbit families.

The final assertion uses ergodicity of the period-one component: an invariant
continuous, hence measurable, charge is constant almost everywhere, and
continuity on regular components gives the stated local conclusion.  Its
increment is therefore zero.  Without invariance, the displayed increment is
by definition \((I-\Ltransfer_0)H\), which is precisely the exact-current
case.

For the finite-memory statement, write \(A(e)\) for the value of \(a\) on an
edge \(e\) of the higher-block graph.  Fix a root vertex \(v_*\).  For every
vertex \(v\), choose a directed path \(P_v\) from \(v_*\) to \(v\), and set
\(K(v)=\sum_{e\in P_v}A(e)\).  This is independent of the chosen path: by
strong connectivity choose a directed return path \(Q_v\) from \(v\) to
\(v_*\); the concatenations \(P_vQ_v\) formed with any two choices of
\(P_v\) are directed cycles, and their \(A\)-sums vanish by the periodic-word
hypothesis.  If \(e:u\to v\), comparison of \(P_u e\) with \(P_v\) gives
\(A(e)=K(v)-K(u)\).  Taking \(v\) to be the current higher-block vertex and
\(u\) its predecessor proves \eqref{eq:finite_memory_label_coboundary}.
\end{proof}

\begin{remark*}[Why the obvious repairs do not produce the missing source]
One may force gluing by projecting a raw source with a right inverse of
\(\mathcal G_{\Gamma,I}\), but this construction uses
\((I-\Ltransfer_0)^{-1}\) and is therefore Poisson-primitive-defined.  A
formal finite-order jet may likewise be corrected recursively with
\(\mathcal R_{\Gamma,K}\), but every correction depends on the preceding
Poisson defect and so is response-tailored rather than independently
prescribed.  A
time-reversal or parity constraint is also insufficient by itself: it removes
the jump only when the symmetry identifies the two one-sided branch trace maps
pointwise, and no such identification is included in the hypotheses of the
nonconjugate radial theorem.
Indeed, suppose a side-swapping involution \(S\) satisfies
\[
 G_s\circ S=-G_s,\qquad
 \mathcal A_{+,s}(h\circ S)
   =(\mathcal A_{-,s}h)\circ S
\]
and the analogous identity with \(+\) and \(-\) exchanged.  If
\(h\circ S=\varepsilon h\), \(\varepsilon\in\{+1,-1\}\), then for
\(J_h=\mathcal A_{+,0}h-\mathcal A_{-,0}h\),
\[
 J_h\circ S=-\varepsilon J_h,\qquad
 \dot G_0\circ S=-\dot G_0,\qquad
 (\dot G_0J_h)\circ S=\varepsilon\dot G_0J_h.
\]
Thus symmetry merely assigns a parity to the face defect; it does not force
that defect to vanish, and an observable of the same parity detects it.
More generally, \eqref{eq:isotypic_gluing_right_inverse} shows that on every
free finite symmetry orbit each isotypic source space still realizes an
arbitrary defect profile on one representative face.  Only a face fixed
pointwise by a stabilizer whose fixed-vector space is trivial could force
the profile to vanish at the abstract level; for finite physical isometries of
the dispersing table,
\cref{prop:no_pointwise_symmetry_stabilizer_grazing_face} rules out that
exception on every regular grazing subarc.
Even a first-order symmetry cancellation would not supply the independent
second source \(q_2\) or its strong remainder.  Thus neither device proves the
nonconjugate prescribed-source theorem left open here.
\end{remark*}
